%! Author = breandan
%! Date = 11/5/25

\documentclass{beamer}

% Basic fonts & symbols
\usepackage[T1]{fontenc}
\usepackage{tgcursor}
\usepackage{booktabs}
\usepackage{pifont}
\newcommand{\cmark}{\ding{51}} % check
\newcommand{\xmark}{\ding{55}} % cross
\usepackage{multicol}
\usepackage[table]{xcolor}

\usepackage{tikz-cd}
\usepackage{adjustbox}

\usepackage{txfonts}
\usepackage{soul}
\usepackage{mathtools}
\usepackage{mathpartir}
\usepackage[os=mac]{menukeys}

\usepackage[normalem]{ulem}
\usepackage{tikz}
\usetikzlibrary{decorations.pathmorphing,calc,automata,arrows}
\usepackage{pgfplots}
\pgfplotsset{compat=1.18}
\usepgfplotslibrary{groupplots}

\usepackage{algorithm}
\usepackage[noend]{algpseudocode}
\usepackage{algorithmicx}

\definecolor{Light}{gray}{.90}
\sethlcolor{Light}
\definecolor{oiBlue}{HTML}{0072B2}
\definecolor{oiGreen}{HTML}{009E73}
\definecolor{oiOrange}{HTML}{E69F00}
\definecolor{oiVermilion}{HTML}{D55E00}
\definecolor{oiBrown}{HTML}{8C564B}

\newcommand{\codett}[1]{\texttt{\hl{#1}}}
\DeclareMathOperator*{\largecomma}{\mathop{\vcenter{\hbox{\Huge\texttt{,}}}}}
\newcommand{\bycommas}{%
  \mathop{\vphantom{\sum}\mathchoice{\Large ,}{,}{,}{,}}\displaylimits
}
\newcommand{\cev}[1]{\reflectbox{\ensuremath{\vec{\reflectbox{\ensuremath{#1}}}}}}
\newcommand\hmmax{0}
\newcommand\bmmax{0}
\usepackage{amssymb}
\usepackage{centernot}
\usepackage{bussproofs}
\DeclareMathAlphabet{\mathcal}{OMS}{cmsy}{m}{n}
\SetMathAlphabet{\mathcal}{bold}{OMS}{cmsy}{b}{n}

\newcommand{\nt}[1]{\texttt{#1}}
\newcommand{\tok}[1]{\codett{#1}}
\newcommand{\gap}{\;\;}
\newcommand{\cjoin}[2]{\mathop{\largecomma}\limits_{#1}^{#2}}

\newcommand{\bs}{\blacksquare}
\newcommand{\ws}{\square}

\newcommand{\T}{\mathbb{T}}
\newcommand{\ttlb}{\{}   % left brace token
\newcommand{\ttrb}{\}}   % right brace token
\newcommand{\ttus}{\_}   % underscore token (for f_*)
\newcommand{\W}{\mathbb W}
\newcommand{\diff}[1]{\textcolor{orange}{#1}}

\pgfplotsset{
  runtimeplot/.style={
    width=0.89\textwidth,
    height=0.58\textheight,
    ymin=0,
    axis line style={draw=black!45, line width=0.35pt},
    tick style={draw=black!45, line width=0.35pt},
    ymajorgrids=true,
    grid style={draw=black!8, line width=0.25pt},
    tick label style={font=\scriptsize},
    label style={font=\scriptsize},
    title style={font=\small},
    every axis plot/.append style={draw=none},
  },
  haltplot/.style={
    width=0.82\textwidth,
    height=0.5\textheight,
    axis line style={draw=black!45, line width=0.35pt},
    tick style={draw=black!45, line width=0.35pt},
    ymajorgrids=true,
    grid style={draw=black!8, line width=0.25pt},
    tick label style={font=\scriptsize},
    label style={font=\scriptsize},
    title style={font=\small},
    every axis plot/.append style={draw=none},
  }
}

% Metadata
\title{Towards a Linear-Algebraic Hypervisor}
\author{Breandan Mark Considine}
\date{\today}

\begin{document}

\begin{frame}
  \titlepage
\end{frame}

\begin{frame}{The RASP Model}
  \vspace{-0.4cm}
  \begin{definition}[Cook \& Reckhow, '73]
  A RASP machine, \(\mathcal{R} = \langle C,c_0,\to,F\rangle\) consists of a configuration space, \( \langle i,a,M,u,y\rangle: C=\mathbb N\times\mathbb Z\times\mathbb Z^{\mathbb N}\times\mathbb Z^*\times\mathbb Z^* \). A RAS program is a vector \(P: (\mathbb Z\times\mathbb Z)^m\) packed with an \textit{input} \(x: \mathbb{Z}^*\), into an \textit{initial configuration}, \(c_0=\big\langle 0,0,M'_{0..2m} \gets P,x0,\varepsilon \big\rangle.\)
  Let \(\langle o,j\rangle=\langle M_i,M_{i+1}\rangle\) denote the \textit{opcode} and \textit{operand} and $\rightarrow$ be defined as:
  {\small
  \[
  \langle i,a,M,u,y\rangle \rightarrow
  \left\{
    \begin{array}{@{}c@{}c@{\:,\:}l@{\:,\:}l@{\:,\:}l@{\:,\:}l@{}c@{\quad}l@{}}
      \langle & i+2 & j     & M               & u & y          & \rangle, & \text{if } o=1_{\mathtt{LOD}},\\
      \langle & i+2 & a+M_j & M               & u & y          & \rangle, & \text{if } o=2_{\mathtt{ADD}},\\
      \langle & i+2 & a-M_j & M               & u & y          & \rangle, & \text{if } o=3_{\mathtt{SUB}},\\
      \langle & i+2 & a     & M_j\gets a      & u & y          & \rangle, & \text{if } o=4_{\mathtt{STO}},\\
      \langle & j   & a     & M               & u & y          & \rangle, & \text{if } o=5_{\mathtt{BPA}}\ \wedge\ a>0,\\
      \langle & i+2 & a     & M               & u & y          & \rangle, & \text{if } o=5_{\mathtt{BPA}}\ \wedge\ a\le 0,\\
      \langle & i+2 & a     & M_j\gets \sigma & v & y          & \rangle, & \text{if } o=6_{\mathtt{RD}}\ \wedge\ u=\sigma v,\\
      \langle & i+2 & a     & M               & u & y\cdot M_j & \rangle, & \text{if } o=7_{\mathtt{PRI}},\\
      \langle & i   & a     & M               & u & y          & \rangle, & \text{otherwise.}
    \end{array}
  \right.
  \]
  }

  and the halting configurations are simply \( F= \{c \in C \mid c \rightarrow c\}\).
    \end{definition}

\end{frame}

\begin{frame}{W-RASP: A Bounded Variant}
  Fix \(w,n,\ell,s\in\mathbb N\), assume \(\W=\mathbb F_2^w\) and define the W-RASP as:
  \[
    \mathcal R^{(w,n,\ell,s)}=\langle C_w,c_0,\to,F\rangle,
    \quad
    C_w=\W\times\W\times\W^n\times\W^{\ell}\times\W^{\sigma}.
  \]
%  \[
%    c_0(P,x)=\big\langle
%    0_w,0_w,P\,0_w^{\,n-2m},x0_w^{\ell-|x|},0_w^{s+1}
%    \big\rangle.
%  \]

  The transition operator takes a configuration $\langle i,a,M,u,y\rangle: C_w$ to:
  \[\small
  \left\{
    \begin{array}{@{}c@{}c@{\:,\:}l@{\:,\:}l@{\:,\:}l@{\:,\:}l@{}c@{\quad}l@{}}
      \langle & i\,\diff{\oplus}\, 2\diff{_w} & j                         & M & u & y                         & \rangle, & \text{if } o=1_{\texttt{LOD}},\\
      \langle & i\,\diff{\oplus}\, 2\diff{_w} & a\,\diff{\oplus}\, M_{j\diff{\bmod n}}       & M & u & y                         & \rangle, & \text{if } o=2_{\texttt{ADD}},\\
      \langle & i\,\diff{\oplus}\, 2\diff{_w} & a\,\diff{\otimes}\, M_{j\diff{\bmod n}}      & M & u & y                         & \rangle, & \text{if } o=3_{\diff{\texttt{MUL}}},\\
      \langle & i\,\diff{\oplus}\, 2\diff{_w} & a                         & M_{j\diff{\bmod n}}\gets a & u & y        & \rangle, & \text{if } o=4_{\texttt{STO}},\\
      \langle & j                             & a                         & M & u & y                         & \rangle, & \text{if } o=5_{\diff{\texttt{BNZ}}}\ \wedge\ a \diff{\neq} 0,\\
      \langle & i\,\diff{\oplus}\, 2\diff{_w} & a                         & M & u & y                         & \rangle, & \text{if } o=5_{\diff{\texttt{BNZ}}}\ \wedge\ a \diff{=} 0,\\
      \langle & i\,\diff{\oplus}\, 2\diff{_w} & a                         & M_{j\diff{\bmod n}}\gets \diff{u^\bullet} & \diff{u^\triangleright} & y & \rangle, & \text{if } o=6_{\texttt{RD}}\ \wedge\ \diff{u_0}\diff{<}\diff{\ell},\\
      \langle & i\,\diff{\oplus}\, 2\diff{_w} & a                         & M & u & y \diff{::} M_{j\diff{\bmod n}}          & \rangle, & \text{if } o=7_{\texttt{PRI}},\\
      \langle & i                             & a                         & M & u & y                         & \rangle, & \text{otherwise.}
    \end{array}
  \right.
  \]
  \[\small
    \langle o,j\rangle:=\big\langle M_{\,i\bmod n},M_{\,(i\,\oplus 1)\bmod n}\big\rangle,
    \qquad
    M_j\gets a := \big[\delta_{jk}a+(1-\delta_{jk})M_k\big]_{k=0}^{n-1}.
  \]
  \[\small
    \diff{u^\bullet}:=u_{\min(u_0+1,\ell)},\qquad
    \diff{u^\triangleright}:=u_0\gets\min(u_0+1,\ell),
  \]
  \[\small
    y\diff{::}z=
    \begin{cases}
      (y_{y_0+1}\gets z)_0\gets y_0+1,&y_0<s,\\
      y,&\text{otherwise.}
    \end{cases}
  \]

\end{frame}

\begin{frame}{Heapless Array Language}
  \vspace{-0.2cm}
  We define a simple array langauge for functions $f: \mathbb{N}^{\lambda\leq \ell} \rightarrow \mathbb{N}^{\eta\leq\sigma}$:
  {
  \[
    \begin{aligned}
      \texttt{F } &::= \texttt{ fun f0 ( ipt : W \char`\^\ N ) -> W \char`\^\ N \{ B \} } \\
      \texttt{N } &::= \texttt{ 1 } \mid \texttt{ 2 }  \mid \texttt{ 3 } \mid \texttt{ 4 } \mid \texttt{ 5 } \mid \texttt{ 6 } \mid \texttt{ 7 } \mid \texttt{ 8 } \mid \texttt{ 9 } \\
      \texttt{B } &::= \texttt{ S } \mid \texttt{ S ; B } \mid \texttt{ hlt } \mid \texttt{ brk} \\
      \texttt{S } &::= \texttt{ P = E } \mid \texttt{ ife G \{ B \} \{ B \} } \mid \texttt{ whl G \{ B \}}\\
      \texttt{Z } &::= \texttt{ 0 } \mid \texttt{ N}\\
      \texttt{G } &::= \texttt{ ipt [ Z ] } \mid \texttt{ scr [ Z ] } \mid \texttt{ opt [ Z ]} \\
      \texttt{P } &::= \texttt{ scr [ Z ] } \mid \texttt{ opt [ Z ]} \\
      \texttt{E } &::= \texttt{ G + G } \mid \texttt{ G * G } \mid \texttt{ G + Z } \mid \texttt{ G * Z } \mid \texttt{ G } \mid \texttt{ Z }\\
    \end{aligned}
  \]
  }

  --\:\:\texttt{ipt}, \texttt{opt} and \texttt{scr} refer to the \textit{input}, \textit{output} and \textit{scratch} buffers.\\--\:\:Grammar grants RO permission for \texttt{ipt} and R+W for \texttt{opt}, \texttt{scr}.\\--\:\:\texttt{hlt} returns immediately and \texttt{brk} escapes the local scope.
\end{frame}

\begin{frame}{Heapless Array Lowering}
  \vspace{-0.7cm}
  \begin{center}
  \setlength{\fboxsep}{2pt}%
  \setlength{\fboxrule}{0.35pt}%
  \fbox{%
  \begin{adjustbox}{max height=0.66\textheight,center}
  \begin{scriptsize}
    \begin{mathpar}
      \inferrule*[right=\scriptsize\textsc{Fun}]
      {\Gamma[\kappa\mapsto \ell_f, \lambda, \eta] \vdash B \Rightarrow C}
      {\vdash \texttt{fun f0 ( ipt : I \char`\^\ }\lambda\texttt{ ) -> I \char `\^\ }\eta\texttt{ \{ }B\texttt{ \}}
      \Rightarrow \keys{\texttt{RD }$\rho_i$\!}\,{}_{i=0}^{\lambda-1}\ C\ {}^{\ell_f:}\keys{\texttt{PRI }$\omega_i$\!}\,{}_{i=0}^{\eta-1}\ \keys{\texttt{HLT}}}

      \inferrule*[right=\scriptsize\textsc{Get}]
      { }
      {\Gamma \vdash \rho[z] \Rightarrow_a \keys{\texttt{LOD }$0$}\ \keys{\texttt{ADD }$\llbracket\rho\rrbracket^\Gamma_z$\!}}
\hspace{0.2cm}
      \inferrule*[right=\scriptsize\textsc{Put}]
      {\Gamma \vdash E \Rightarrow_a A}
      {\Gamma \vdash \rho[z] \texttt{ = } E \Rightarrow A\ \keys{\texttt{STO }$\llbracket\rho\rrbracket^\Gamma_z$\!}}

      \inferrule*[right=\scriptsize\(\oplus\)]
      {\Gamma \vdash I_1 \Rightarrow_a J_1 \\
      \Gamma \vdash I_2 \Rightarrow_a J_2}
      {\Gamma \vdash I_1 + I_2 \Rightarrow_a J_1\ \keys{\texttt{STO }$\nu$}\ J_2\ \keys{\texttt{ADD }$\nu$}}
      \hspace{0.1cm}
      \inferrule*[right=\scriptsize\(\otimes\)]
      {\Gamma \vdash I_1 \Rightarrow_a J_1 \\
      \Gamma \vdash I_2 \Rightarrow_a J_2}
      {\Gamma \vdash I_1 \times I_2 \Rightarrow_a J_1\ \keys{\texttt{STO }$\nu$}\ J_2\ \keys{\texttt{MUL }$\nu$}}

      \inferrule*[right=\scriptsize\textsc{Ife}]
      {\Gamma \vdash H \Rightarrow_a J \\
      \Gamma \vdash B_1 \Rightarrow C_1 \\
      \Gamma \vdash B_2 \Rightarrow C_2}
      {\Gamma \vdash \texttt{ife}\ H\ \texttt{\{}\ B_1\ \texttt{\}}\ \texttt{\{}\ B_2\ \texttt{\}}
      \Rightarrow
      J\ \keys{\texttt{BNZ }$\ell_t$\!}\ C_2\ \keys{\texttt{JMP }$\ell_e$\!}\ {}^{\ell_t:}\,C_1\ {}^{\ell_e:}}

      \inferrule*[right=\scriptsize\textsc{While}]
      {\Gamma \vdash H \Rightarrow_a J \\
      \Gamma[\kappa\mapsto \ell_e] \vdash B \Rightarrow C}
      {\Gamma \vdash \texttt{whl}\ H\ \texttt{\{}\ B\ \texttt{\}}
      \Rightarrow
      {}^{\ell_h:}\,J\ \keys{\texttt{BNZ }$\ell_b$\!}\ \keys{\texttt{JMP }$\ell_e$\!}\
      {}^{\ell_b:}\,C\ \keys{\texttt{JMP }$\ell_h$\!}\ {}^{\ell_e:}}
      \hspace{0.1cm}
      \inferrule*[right=\scriptsize\textsc{Break}]
      {\Gamma(\kappa)=\ell}
      {\Gamma \vdash \texttt{brk} \Rightarrow \keys{\texttt{JMP }$\ell$}}

      \inferrule*[right=\scriptsize\textsc{Seq}]
      {\Gamma \vdash B \Rightarrow C_B \\
      \Gamma \vdash S \Rightarrow C_S}
      {\Gamma \vdash B\texttt{ ; }S \Rightarrow
      C_B\ C_S}
      \hspace{0.3cm}
      \inferrule*[right=\scriptsize\textsc{Halt}]
      { }
      {\Gamma \vdash \texttt{hlt} \Rightarrow
      \keys{\texttt{PRI }$\omega_i$\!}\,{}_{i=0}^{l-1}\ \keys{\texttt{HLT}}}
    \end{mathpar}
  \end{scriptsize}
  \end{adjustbox}%
  }
  \end{center}

  \vspace{-0.1cm}
  \small We write \underline{\(\Gamma \vdash E \Rightarrow A\)} when the \(A\) implements \(E\), and \underline{\(\Gamma \vdash E \Rightarrow_a A\)} when \(a\) is left in the accumulator. Desugar \keys{\texttt{JMP }$\ell$} as \keys{\texttt{LOD }$1$} \keys{\texttt{BNZ }$\ell$} and interpret array indexing modulo declared length, i.e., \( \llbracket\rho\rrbracket^\Gamma_z := M_{2m + \llbracket\rho\rrbracket + (z\bmod |\rho|)}. \)
\end{frame}

\begin{frame}{Memory Layout}
  \small
  \vspace{-0.25cm}
  Reserve disjoint memory regions \(\alpha_{0\dots \lambda-1}\), \(\omega_{0\dots \eta-1}\), and \(\zeta_{0\dots \gamma-1}\) for arrays \(\texttt{ipt}: \mathbb{W}^{\lambda \le \ell}\), \(\texttt{opt}: \mathbb{W}^{\eta \le \sigma}\) and \(\texttt{scr}: \mathbb{W}^{\gamma \le \mu}\), with \((5+\ell+\sigma+\mu+2m)\le n\). The hypervisor stores each VM as one contiguous 32-bit row vector. RASP instructions address only \(M_v\).

  \vspace{0.4cm}

  \begin{adjustbox}{max width=\textwidth}
  \(
  M_v =
  \big[
    \overbrace{
      \underset{\tiny\keys{\texttt{LOD}\ 0}}{1},
      \underset{\tiny\keys{\texttt{ADD}\ $\alpha_{\scalebox{0.65}{$\scriptscriptstyle 0$}}$\!}}{66},
      \underset{\tiny\keys{\texttt{STO}\ $\omega_{\scalebox{0.65}{$\scriptscriptstyle 0$}}$\!}}{76},
      \underset{\tiny\keys{\texttt{LOD}\ 1}}{9},
      \underset{\tiny\keys{\texttt{BNZ}\ 6}}{53},
      \underset{\tiny\keys{\texttt{PRI}\ $\omega_{\scalebox{0.65}{$\scriptscriptstyle 0$}}$\!}}{79},
      \underset{\tiny\keys{\texttt{HLT}}}{0},\underbrace{0,\ldots,0}_{\text{padding}}
    }^{P_{0..c-1}:\ \text{packed code}},
    \underset{\lambda}{3},\overbrace{\underbrace{0,\ldots}_{|\alpha|=\lambda},0}^{\mathtt{ipt}=\alpha},
    \underset{\eta}{2},\overbrace{\underbrace{0,\ldots}_{|\omega|=\eta},0}^{\mathtt{opt}=\omega},
    \underset{\gamma}{4},\overbrace{\underbrace{0,\ldots}_{|\zeta|=\gamma},0}^{\mathtt{scr}=\zeta},
    \underset{\nu}{0},
    \overbrace{0,\ldots,0}^{\text{padding}}
  \big]
  \)
  \end{adjustbox}

  \vspace{.5cm}

  \begin{adjustbox}{max width=\textwidth}
  \(
  VM_v =
  \overbrace{[
  \mathtt{pc},\mathtt{acc},\mathtt{steps},\mathtt{halted},
  \mathtt{outCount},\mathtt{out}_0,\mathtt{out}_1,\mathtt{out}_2
  ]}^{H_v:\ \text{hypervisor sidecar}}
  \ \Vert\
  \overbrace{\big[M_v[0],\ldots,M_v[n-1]\big]}^{M_v:\ \text{RASP-addressable}},
  \quad |VM_v| \approx 2^8.
  \)
  \end{adjustbox}

  \vspace{.5cm}

  \begin{adjustbox}{max width=\textwidth}
  \(
  HVM = [\overbrace{VM_1, VM_2, \underline{VM_3}, \ldots}^{t_1}, \overbrace{VM_1, VM_2, \underline{VM_3}, \ldots}^{t_2}, \ldots, \overbrace{VM_1, VM_2, \underline{VM_3}, \ldots}^{t_{W-1}}]
  \)
  \end{adjustbox}

  \vspace{.5cm}

Lift the transition function to a vectorized transition map, \(\Phi:C^d\to C^d\). The resulting configuration after \(i\) steps is \( \mathbf c_i=\Phi^i(\mathbf c_0).\) For a vector of programs \(\mathbf{P}\) and inputs \(\mathbf x\), we will write \( \mathbf c_i(\mathbf P,\mathbf x):=\Phi^i\bigl(\mathbf c_0(\mathbf P,\mathbf x)\bigr)\) for the evolution of $n$ independent machines, each with static memory.

\end{frame}

\begin{frame}{RASP Hypervisor}
\begin{algorithm}[H]
\caption{Striped round-robin RASP hypervisor}
\label{alg:hypervisor}
\small
\begin{algorithmic}[1]
\State \textbf{inputs:} Global config \(\mathbf c_0 \in C^d\), epoch \(q\), step budget \(\tau_{\max}\)
\State \(\mathbf c \gets \mathbf c_0\), \(W \gets \textsc{MaxResidentThreads}()\)
\For{\(g = 0\) \textbf{to} \(W-1\) \textbf{in parallel}} \Comment{Persistent workers}
\For{\(r = 0\) \textbf{to} \(\lceil \tau_{\max} / q \rceil - 1\)} \Comment{Total rounds}
\For{\(k = 0\) \textbf{to} \(\lceil d / W \rceil-1\)} \Comment{Walk worker stripe}
\State \(j \gets g + kW\) \Comment{Current VM index}
\If{\(j < d\) \textbf{and} \(\mathbf c[j] \notin F\)} \Comment{Skip halted VMs}
\State \(R \gets \textsc{LoadRegs}(\mathbf c[j])\)
\For{\(t = 1\) \textbf{to} \(q\)} \Comment{Bounded epoch}
\If{\(R \in F\)} \textbf{break} \EndIf
\State \(R \gets \Phi_w(R)\) \Comment{One-step transition}
\EndFor
\State \(\mathbf c[j] \gets \textsc{StoreRegs}(R)\)
\EndIf
\EndFor
\EndFor
\EndFor
\State \textbf{return} \(\mathbf c\)
\end{algorithmic}
\end{algorithm}
\end{frame}

\begin{frame}{Evaluation: Runtime}
  \scriptsize
  \centering
  \begin{tikzpicture}
    \begin{axis}[
      runtimeplot,
      title={Average runtime},
      ylabel={Time (s)},
      xlabel={Total VMs $(\times 10^6)$},
      ymode=log,
      log basis y={10},
      ymin=0.5,
      ymax=500,
      xtick={1,2,3,4,5,6,7,8},
      ybar=0.4pt,
      bar width=2.4pt,
      legend image post style={scale=0.6},
      legend style={font=\fontsize{5.2}{6}\selectfont, at={(0.01,0.99)}, anchor=north west, draw=none, fill=none, legend columns=3}
    ]
      \addplot+[fill=gray] coordinates {(1,46.5) (2,95.2) (3,143.8) (4,188.9) (5,239.1) (6,282.0) (7,327.1) (8,382.6)};
      \addlegendentry{M4 ($\times1$)}
      \addplot+[fill=oiBlue] coordinates {(1,4.6) (2,8.9) (3,13.1) (4,18.6) (5,22.2) (6,26.8) (7,30.8) (8,35.3)};
      \addlegendentry{M4 ($\times16$)}
      \addplot+[fill=oiGreen] coordinates {(1,2.2) (2,4.3) (3,6.3) (4,8.3) (5,10.4) (6,12.4) (7,14.4) (8,16.5)};
      \addlegendentry{A10G}
      \addplot+[fill=oiBrown] coordinates {(1,1.1) (2,1.9) (3,2.8) (4,3.6) (5,4.5) (6,5.3) (7,6.2) (8,7.0)};
      \addlegendentry{B200 ($\times1$)}
      \addplot+[fill=oiVermilion] coordinates {(1,0.6) (2,0.9) (3,1.2) (4,1.5) (5,1.8) (6,2.0) (7,2.3) (8,2.6)};
      \addlegendentry{B200 ($\times8$)}
    \end{axis}
  \end{tikzpicture}

  \vspace{0.15cm}
  \raggedright
  \(w=32\), \(n=250\), \(\ell=10\), \(s=2\), \(\tau_{\max}=10^6\), \(d=8\times10^6\), \(L=100\).
\end{frame}

\begin{frame}{Evaluation: Halting Distribution}

  \[
    \Pr_{P' \sim \mathcal P_L}\Big(\Phi_w^{\tau_{\max}}\bigl(c_0(P'\Rightarrow P, x\sim \mathbb{F}_2^{w\cdot\text{arity}(P')})\bigr)\in F \bigm\vert L = |P'|\Big).
  \]
  \centering
  \begin{tikzpicture}
    \begin{axis}[
      haltplot,
      title={Halting distribution},
      ylabel={Count},
      xlabel={Halting steps $(\tau_h)$},
      ybar,
      ymode=log,
      yticklabels={0,$10^2$,$10^4$,$10^6$},
      log ticks with fixed point,
      minor y tick num=4,
      yticklabel style={/pgf/number format/fixed},
      bar width=1pt,
      ymin=1,
      ymax=1000000,
      xmin=0,
      xmax=103,
      enlarge x limits=false,
      xtick={10,20,30,40,50,60,70,80,90,100,101},
      xticklabels={10,20,30,40,50,60,70,80,90,{$[10^2,10^6)$},{$10^6{\sim}\infty$}},
      xticklabel style={rotate=45, anchor=east},
    ]
      \addplot+[fill=lightgray] coordinates {
        (1,0) (2,0) (3,0) (4,0) (5,11949) (6,11828) (7,72097) (8,75214)
        (9,115990) (10,131440) (11,156142) (12,182530) (13,215934) (14,232771)
        (15,264785) (16,231876) (17,256198) (18,242506) (19,257656) (20,261201)
        (21,261929) (22,255723) (23,252844) (24,246236) (25,243821) (26,241862)
        (27,235431) (28,229969) (29,221729) (30,214793) (31,208180) (32,200879)
        (33,194560) (34,187593) (35,178581) (36,169924) (37,162759) (38,153716)
        (39,145518) (40,135943) (41,126372) (42,116195) (43,108277) (44,97265)
        (45,87474) (46,76340) (47,65987) (48,56914) (49,47056) (50,38013)
        (51,29896) (52,21811) (53,16876) (54,11536) (55,7780) (56,5730)
        (57,2643) (58,1894) (59,744) (60,391) (61,322) (62,120) (63,61)
        (64,52) (65,30) (66,23) (67,21) (68,17) (69,8) (70,12) (71,8)
        (72,8) (73,6) (74,6) (75,3) (76,2) (77,3) (78,1) (79,4) (80,7)
        (81,1) (82,5) (83,1) (84,5) (85,0) (86,1) (87,2) (88,1) (89,4)
        (90,1) (91,0) (92,1) (93,0) (94,1) (95,1) (96,1) (97,3) (98,0) (99,0)
      };
      \addplot+[fill=oiOrange, draw=oiOrange, bar width=0.7pt] coordinates { (100,664) };
      \addplot+[fill=oiVermilion, draw=oiVermilion, bar width=0.7pt] coordinates { (101,217293) };
      \node[oiOrange, align=right, anchor=east, font=] at (axis cs:100, 664) {\fontsize{4.6}{4}\selectfont$[10^2,10^6)\to$};
      \node[oiVermilion, align=right, anchor=east, font=] at (axis cs:103, 220000) {\fontsize{4.6}{4}\selectfont$10^6\le \tau_h \approx \infty\to$};
    \end{axis}
  \end{tikzpicture}

  \vspace{0.05cm}
  \raggedright
  We plot programs by their halting number \(\tau_h\), the least \(\tau\) such that \(\Phi_w^\tau(c_0)=\Phi_w^{\tau+1}(c_0)\).
\end{frame}

\end{document}
